% GENERATED FILE. Edit canonical lessons, then run npm run generate:books.
% canonical-source-hash: fnv1a64:5f782de2ca069541
% canonical-lessons: JA-C165-nanoka, JA-C165-youka, JA-C165-tooka, JA-C165-hatsuka, JA-C165-tanjoubi

\chapter{The Seventh, the Tenth and a Birthday}
\label{ch:ja-the-seventh-the-tenth-and-a-birthday}
\hlchaptermodality{165}

\begin{chapteropening}
\textbf{By the end of this chapter:} \emph{I can say dates in the month and plan for a birthday.}

Complete the last lesson of chapter 165: I can say dates in the month and plan for a birthday.
\end{chapteropening}

\section[\texorpdfstring{\ja{なのか} (nanoka)}{なのか (nanoka)}]{\ja{なのか} (nanoka) — the seventh of the month; seven days}

\label{lesson:JA-C165-nanoka}



\begin{quote}
Before the new one: say the Japanese for Friday, then the Japanese for Sunday.
\end{quote}

\subsection*{You'll want to know: \ja{なのか}}
\textbf{\ja{なのか}} — \emph{nanoka} — \textquotedblleft{}the seventh of the month; seven days\textquotedblright{}.

\begin{quote}
\textbf{\ja{らいげつの} \ja{なのかに} \ja{うみに} \ja{いきます。}} — \emph{raigetsu no nanoka ni umi ni ikimasu} — \textquotedblleft{}I will go to the sea on the seventh of next month.\textquotedblright{}
\end{quote}

\begin{cousinweb}
\textbf{\ja{なのか}} (\emph{nanoka}) is native Japanese, built from \emph{nano-}, a form of \emph{nana}, \textquotedblleft{}seven\textquotedblright{}, and \emph{ka}, \textquotedblleft{}day\textquotedblright{}, both inherited rather than borrowed (JMdict, KANJIDIC2).
\end{cousinweb}

\begin{grammarlens}[title={What you've built}]
One more word for everyday life at A2.
\end{grammarlens}

\subsection*{Your turn}
\begin{itemize}
\raggedright
  \item \emph{Say it:} \emph{nanoka}
  \item \emph{Say it:} \emph{nanoka}, once more
  \item \emph{Say it:} \emph{raigetsu no nanoka ni umi ni ikimasu} — I will go to the sea on the seventh of next month.
  \item \emph{Recall:} say the Japanese for Friday, then the Japanese for Sunday, then say \emph{nanoka} again
\end{itemize}

\begin{tcolorbox}[breakable,colback=teal!4,colframe=teal!35!black,title={Before you move on}]
What does \ja{なのか} mean? (\textquotedblleft{}The seventh of the month; seven days\textquotedblright{}.) Say it once more.
\end{tcolorbox}
\section[\texorpdfstring{\ja{ようか} (yōka)}{ようか (yōka)}]{\ja{ようか} (yōka) — the eighth of the month; eight days}

\label{lesson:JA-C165-youka}



\begin{quote}
Before the new one: say the Japanese for Sunday, then the Japanese for the seventh of the month.
\end{quote}

\subsection*{You'll want to know: \ja{ようか}}
\textbf{\ja{ようか}} — \emph{yōka} — \textquotedblleft{}the eighth of the month; eight days\textquotedblright{}.

\begin{quote}
\textbf{\ja{ようかは} \ja{にちようびでしょう。}} — \emph{yōka wa nichiyōbi deshō} — \textquotedblleft{}The eighth is probably a Sunday.\textquotedblright{}
\end{quote}

\begin{cousinweb}
\textbf{\ja{ようか}} (\emph{yōka}) is native Japanese, built from \emph{ya}, \textquotedblleft{}eight\textquotedblright{}, which takes the form \emph{yō} here, and \emph{ka}, \textquotedblleft{}day\textquotedblright{}, both inherited rather than borrowed (JMdict, KANJIDIC2).
\end{cousinweb}

\begin{grammarlens}[title={What you've built}]
One more word for everyday life at A2.
\end{grammarlens}

\subsection*{Your turn}
\begin{itemize}
\raggedright
  \item \emph{Say it:} \emph{yōka}
  \item \emph{Say it:} \emph{yōka}, once more
  \item \emph{Say it:} \emph{yōka wa nichiyōbi deshō} — The eighth is probably a Sunday.
  \item \emph{Recall:} say the Japanese for Sunday, then the Japanese for the seventh of the month, then say \emph{yōka} again
\end{itemize}

\begin{tcolorbox}[breakable,colback=teal!4,colframe=teal!35!black,title={Before you move on}]
What does \ja{ようか} mean? (\textquotedblleft{}The eighth of the month; eight days\textquotedblright{}.) Say it once more.
\end{tcolorbox}
\section[\texorpdfstring{\ja{とおか} (tōka)}{とおか (tōka)}]{\ja{とおか} (tōka) — the tenth of the month; ten days}

\label{lesson:JA-C165-tooka}



\begin{quote}
Before the new one: say the Japanese for the seventh of the month, then the Japanese for the eighth of the month.
\end{quote}

\subsection*{You'll want to know: \ja{とおか}}
\textbf{\ja{とおか}} — \emph{tōka} — \textquotedblleft{}the tenth of the month; ten days\textquotedblright{}.

\begin{quote}
\textbf{\ja{とおかに} \ja{ともだちと} \ja{あう} \ja{つもりです。}} — \emph{tōka ni tomodachi to au tsumori desu} — \textquotedblleft{}I plan to meet a friend on the tenth.\textquotedblright{}
\end{quote}

\begin{cousinweb}
\textbf{\ja{とおか}} (\emph{tōka}) is native Japanese, built from \emph{tō}, \textquotedblleft{}ten\textquotedblright{}, and \emph{ka}, \textquotedblleft{}day\textquotedblright{}, both inherited rather than borrowed (JMdict, KANJIDIC2).
\end{cousinweb}

\begin{grammarlens}[title={What you've built}]
One more word for everyday life at A2.
\end{grammarlens}

\subsection*{Your turn}
\begin{itemize}
\raggedright
  \item \emph{Say it:} \emph{tōka}
  \item \emph{Say it:} \emph{tōka}, once more
  \item \emph{Say it:} \emph{tōka ni tomodachi to au tsumori desu} — I plan to meet a friend on the tenth.
  \item \emph{Recall:} say the Japanese for the seventh of the month, then the Japanese for the eighth of the month, then say \emph{tōka} again
\end{itemize}

\begin{tcolorbox}[breakable,colback=teal!4,colframe=teal!35!black,title={Before you move on}]
What does \ja{とおか} mean? (\textquotedblleft{}The tenth of the month; ten days\textquotedblright{}.) Say it once more.
\end{tcolorbox}
\section[\texorpdfstring{\ja{はつか} (hatsuka)}{はつか (hatsuka)}]{\ja{はつか} (hatsuka) — the twentieth of the month; twenty days}

\label{lesson:JA-C165-hatsuka}



\begin{quote}
Before the new one: say the Japanese for the eighth of the month, then the Japanese for the tenth of the month.
\end{quote}

\subsection*{You'll want to know: \ja{はつか}}
\textbf{\ja{はつか}} — \emph{hatsuka} — \textquotedblleft{}the twentieth of the month; twenty days\textquotedblright{}.

\begin{quote}
\textbf{\ja{はつかに} \ja{いなかに} \ja{かえります。}} — \emph{hatsuka ni inaka ni kaerimasu} — \textquotedblleft{}I will go back to the countryside on the twentieth.\textquotedblright{}
\end{quote}

\begin{cousinweb}
\textbf{\ja{はつか}} (\emph{hatsuka}) is native Japanese, built from an old native word for \textquotedblleft{}twenty\textquotedblright{} and \emph{ka}, \textquotedblleft{}day\textquotedblright{}, both inherited rather than borrowed; the older origin of the first part is not settled (JMdict, KANJIDIC2).
\end{cousinweb}

\begin{grammarlens}[title={What you've built}]
One more word for everyday life at A2.
\end{grammarlens}

\subsection*{Your turn}
\begin{itemize}
\raggedright
  \item \emph{Say it:} \emph{hatsuka}
  \item \emph{Say it:} \emph{hatsuka}, once more
  \item \emph{Say it:} \emph{hatsuka ni inaka ni kaerimasu} — I will go back to the countryside on the twentieth.
  \item \emph{Recall:} say the Japanese for the eighth of the month, then the Japanese for the tenth of the month, then say \emph{hatsuka} again
\end{itemize}

\begin{tcolorbox}[breakable,colback=teal!4,colframe=teal!35!black,title={Before you move on}]
What does \ja{はつか} mean? (\textquotedblleft{}The twentieth of the month; twenty days\textquotedblright{}.) Say it once more.
\end{tcolorbox}
\section[\texorpdfstring{\ja{たんじょうび} (tanjōbi)}{たんじょうび (tanjōbi)}]{\ja{たんじょうび} (tanjōbi) — a birthday}

\label{lesson:JA-C165-tanjoubi}



\begin{quote}
Before the new one: say the Japanese for the tenth of the month, then the Japanese for the twentieth of the month.
\end{quote}

\subsection*{You'll want to know: \ja{たんじょうび}}
\textbf{\ja{たんじょうび}} — \emph{tanjōbi} — \textquotedblleft{}a birthday\textquotedblright{}.

\begin{quote}
\textbf{\ja{あしたは} \ja{いもうとの} \ja{たんじょうびです。}} — \emph{ashita wa imōto no tanjōbi desu} — \textquotedblleft{}Tomorrow is my younger sister's birthday.\textquotedblright{}
\end{quote}

\begin{cousinweb}
\textbf{\ja{たんじょうび}} (\emph{tanjōbi}) is a mixed word: Sino-Japanese \emph{tan}, \textquotedblleft{}birth\textquotedblright{} (Middle Chinese \emph{danX}), and \emph{jō}, \textquotedblleft{}life, be born\textquotedblright{} (\emph{srjaeng}), borrowed from Chinese, with native \emph{hi}, \textquotedblleft{}day\textquotedblright{}, voiced to \emph{bi} (JMdict, KANJIDIC2; Middle Chinese in Baxter's transcription).
\end{cousinweb}

\begin{grammarlens}[title={What you've built}]
One more word for everyday life at A2.
\end{grammarlens}

\subsection*{Your turn}
\begin{itemize}
\raggedright
  \item \emph{Say it:} \emph{tanjōbi}
  \item \emph{Say it:} \emph{tanjōbi}, once more
  \item \emph{Say it:} \emph{ashita wa imōto no tanjōbi desu} — Tomorrow is my younger sister's birthday.
  \item \emph{Recall:} say the Japanese for the tenth of the month, then the Japanese for the twentieth of the month, then say \emph{tanjōbi} again
\end{itemize}

\begin{tcolorbox}[breakable,colback=teal!4,colframe=teal!35!black,title={Before you move on}]
What does \ja{たんじょうび} mean? (\textquotedblleft{}A birthday\textquotedblright{}.) Say it once more.
\end{tcolorbox}
